% !TeX TS-program = pdflatex
\documentclass[10pt,a4paper]{article}

\usepackage[utf8]{inputenc}
\usepackage[T1]{fontenc}
\usepackage[a4paper, margin=1.1cm]{geometry}
\usepackage{enumitem}
\usepackage{hyperref}
\usepackage{microtype}
\usepackage{titlesec}
\usepackage{lmodern,textcomp}

\hypersetup{
    pdfauthor={Richard Harnisch},
    pdftitle={Richard Harnisch - Curriculum Vitae},
    pdfsubject={Curriculum Vitae},
    pdfkeywords={Richard Harnisch, CV, Machine Learning},
    pdfproducer={LaTeX},
    pdfcreator={pdflatex}
}

\pagestyle{empty}
\setlength{\parindent}{0pt}
\setlength{\parskip}{0.09em}
\setlist[itemize]{leftmargin=1.15em, itemsep=0.05em, topsep=0.12em, parsep=0pt, partopsep=0pt}
\titleformat{\section}{\bfseries\large}{}{0pt}{}[\titlerule]
\titlespacing*{\section}{0pt}{0.48em}{0.1em}
\newcommand{\cvheading}[4]{\textbf{#1}\hfill #2\\\emph{#3}\hfill #4\par}

\begin{document}

\begin{center}
    {\LARGE \textbf{Richard Harnisch}}\\[-0.05em]
    \href{mailto:ric.harnisch@gmail.com}{ric.harnisch@gmail.com} \quad
    \href{https://www.linkedin.com/in/richardharnisch}{linkedin.com/in/richardharnisch} \quad
    \href{https://github.com/richardharnisch}{github.com/richardharnisch}
\end{center}

\section*{Education}
\cvheading{MSc Artificial Intelligence, current GPA 9.0}{September 2026--2028}{University of Amsterdam}{}
\begin{itemize}
    \item Focusing on Reinforcement Learning and LLMs, working on extracurricular project with \textbf{Pangram} on \textbf{LLM detection}.
\end{itemize}
\cvheading{BSc Artificial Intelligence, cum laude, GPA 8.4/10}{2022--2026}{University of Groningen}{}
\begin{itemize}
    \item Thesis: Interpolation Threshold and \textbf{Double Descent in RL}. Tracked Fisher Information Matrix trace across over-parameterized agents; found local curvature can decouple from the generalization gap under policy-induced non-stationarity.
    \item Formal math/CS foundations: Calculus for Artificial Intelligence; Linear Algebra and Multivariable Calculus; Statistics; Analysis; Probability Theory; Probability and Measure; Linear Algebra 2; Partial Differential Equations; Imperative Programming; Algorithms and Data Structures; Object-Oriented Programming.
    \item Advanced AI/ML: Applied Machine Learning, \textbf{Reinforcement Learning}, Uncertainty in ML, Machine Translation (\textbf{9.5 each}); Neural Networks, Natural Language Processing (9 each); Advanced Logic (8.5).
\end{itemize}
\cvheading{International Baccalaureate; German Abitur 1.0}{2020--2022}{Dresden International School}{42/45 points}
Higher Level Mathematics, Physics, Economics; \textbf{maximum marks} in Mathematics, Physics, German, Spanish. Extended Essay on public-key cryptography including Euler's totient function, Chinese remainder theorem \& vulnerabilities.

\section*{Research and ML Experience}
\cvheading{Research Assistant: Data Center Optimization}{2026--present}{Multiscale Networked Systems Group}{Amsterdam}
\begin{itemize}
    \item Working at the University of Amsterdam under Dr. Shashikant Ilager on Data Center design, targeting publication.
    \item Optimizing electricity use and carbon intensity while keeping performance within threshold values for both inference and training workloads. Exploring \textbf{agent-driven design} for physically verified data-center designs while minimizing cost.
\end{itemize}
\cvheading{Data Science Intern: Electricity Price Forecasting}{2025--2026}{Nederlandse Gasunie}{Groningen/Remote}
\begin{itemize}
    \item Developed a day-ahead electricity price forecasting system using SARIMA, LSTMs, and Transformers in 5-person team. % under the Chief Data Scientist.
    \item Designed and deployed \textbf{distribution-shift monitoring} over ENTSO-E price data and Dutch grid-production data; system alerts operators and triggers retraining under covariate drift; worked \textbf{directly with Chief Data Scientist}.
    \item Conducted feature engineering and data management for \textbf{real-world forecasting} in a grid undergoing renewable- and geopolitics-driven distribution shifts including alerting operator upon distribution shift to retrain model on new data.
\end{itemize}
\cvheading{Bachelor's Thesis: Double Descent in Reinforcement Learning}{2025--2026}{University of Groningen}{Supervisor: Dr.\ Matthia Sabatelli}
\begin{itemize}
    \item Investigated the \textbf{interpolation threshold} and Double Descent phenomenon in \textbf{over-parameterized Reinforcement Learning agents} under non-stationary data distributions created by \textbf{randomized environments} and adaptive policy.
    \item Tracked the Fisher Information Matrix (FIM) trace as a \textbf{generalization metric} across constrained policy optimizers (TRPO) at varying model capacity and training durations. Gained experience with \textbf{high-performance computing}.
    \item Found empirical evidence that FIM decorrelates from the generalization gap under policy-induced non-stationarity: constrained optimizers exhibit expected FIM spike near the interpolation threshold, but the curvature signal does not predict generalization gap evolution, challenging the applicability of classical information-geometric bounds in RL.
    \item Raises open questions about the role of implicit regularization and optimizer geometry in non-stationary generalization.
\end{itemize}
\cvheading{Machine Learning Engineer and Software Developer}{2024--2026}{descript GmbH}{Dresden/Remote}
\begin{itemize}
    \item Independently \textbf{led industrial ML research collaboration with TU Dresden} on textile defect detection; defined project methodology for learning from limited data and built \textbf{active learning} system for prioritized expert review.
    \item Led \textbf{full SWE pipeline for production features} from issue to deployment, such as a building management scheduler.
\end{itemize}
\cvheading{Teaching Assistant: Introduction to Machine Learning}{2025--2026}{University of Groningen}{}
\begin{itemize}
    \item Led lab sessions and graded assignments for 100 second-year AI students, communicated ML concepts to broad audience.
    \item Topics: ML paradigms, model selection, regularization, evaluation metrics, data ingestion, processing and augmentation.
\end{itemize}
% \cvheading{Commissioner of Educational Affairs}{2023--2024}{Study Association Cover}{Groningen}
% \begin{itemize}
%     \item Board of association with 160k€ annual spending, organized extracurricular ML lecture series in collaboration with faculty; designed and led RL chess-bot competition including competition infrastructure; Board of Advisors 2024--2025.
% \end{itemize}
% \cvheading{Software Engineering Intern}{2019}{descript GmbH}{Dresden}
% \begin{itemize}
%     \item Designed and deployed a live CI/CD pipeline status monitoring system using GitLab API and Vue.js.
% \end{itemize}

\section*{Selected Projects}
\textbf{Hackathon Winner.} With two teammates, won AI Hackathon Haarlem 2026 with \textbf{agentic orchestration system} to ingest house style based on example \texttt{pptx} files and delivers arbitrary editable \texttt{pptx} decks including speaker notes \& handouts.\\
\textbf{Study Association Board.} As board of \textbf{1,000+ member AI/CS study association} at University of Groningen \textbf{organized AI symposium} with academic \& industry speakers; \textasciitilde 150 attendees. Organized international excursions.\\
\textbf{Mars Rover Competition.} \textbf{Authored technical report} for European Space Agency rover competition and presented with MakerCie team; placed \textbf{1st among 33 teams} and received maximum points for presentation.\\
\textbf{Open Source Contributions.} Submitted bug-fix PRs to both \textbf{Opencode} and \textbf{PyTorch (vision)}, accepted and merged.
\vfill
\textbf{Languages:} \textbf{English} (TOEFL iBT 6/6, C2) and \textbf{German} native; \textbf{Dutch} and \textbf{Spanish} conversational.

\end{document}
